\chapter{2. Ricci curvature bounded below}
\source{cheeger-gromov-taylor:page-0010,cheeger-gromov-taylor:page-0011,cheeger-gromov-taylor:page-0012,cheeger-gromov-taylor:page-0013,cheeger-gromov-taylor:page-0014,cheeger-gromov-taylor:page-0015,cheeger-gromov-taylor:page-0016,cheeger-gromov-taylor:page-0017,cheeger-gromov-taylor:page-0018,cheeger-gromov-taylor:page-0019}

We point out that Theorem 1.4 can be extended to the case $\partial M^n \neq \emptyset$, by
use of a suitable generalization of Proposition 4.2.

\bigskip

Let $M^n$ be an arbitrary complete Riemannian manifold with $\partial M^n = \emptyset$.
Let $\Tscr_2^{\phi,A}$ be the class of even functions $\{f(\lambda)\}$, such that for some constants
$A$, $k_0$, $s_0$ we have

\begin{equation}\tag{2.1}
\bigl|f^{(k)}(s)\bigr| \le c_1\left(\frac{k}{A}\right)^k \phi(s),\qquad k \ge k_0,\ s \ge s_0,
\end{equation}

where $\phi \ge 0$ belongs to $L^1(a,\infty)$ for all $a > 0$. If we set

\begin{equation}\tag{2.2}
\psi(r) = \int_r^\infty \phi(s)\,ds,
\end{equation}

then

\begin{equation}\tag{2.3}
\int_r^\infty \bigl|f^{(k)}(s)\bigr|\,ds \le c_1\left(\frac{k}{A}\right)^k \psi(r).
\end{equation}

We are going to estimate $k_f$ for $f \in \Tscr_2^{\phi,A}$. Our first step is to reduce the
problem to estimating solutions to the homogeneous Laplace equation
(harmonic functions) in dimension $(n+1)$. The device we employ was used by
Taylor [34], and Lions and Magenes [27], in their independently discovered
proof of the Kotake-Narasimhan Theorem. Observe that if $\rho(x_1, x_2) = d =
2A + r$, by Corollary 1.2 we have

\begin{equation}\tag{2.4}
\bigl\|\Delta^k f(\sqrt{-\Delta})\Delta^l u\bigr\|_{B_A(x_1)}
\le c\left(\frac{2k + 2l}{A}\right)^{2k+2l}\psi(r)\|u\|_{B_A(x_2)},
\end{equation}

for any $u \in L^2$ with $\supp u \subset B_r(x_2)$. Thus for such $u$, for $|y| \le A/2$, the
power series

\begin{equation}\tag{2.5}
\begin{aligned}
v(x,y) &= \sum_{k=0}^\infty \frac{y^{2k}}{(2k)!}(-\Delta)^k f(\sqrt{-\Delta})u,\\
&= (\cosh \sqrt{-\Delta}\, y)\,f(\sqrt{-\Delta})u(x),\ (\mathrm{formally})
\end{aligned}
\end{equation}

converges to an element of $L^2([ -A/2,\, A/2]\times B_r(x_1))$. Moreover,

\begin{equation}\tag{2.6}
\left(\Delta_M + \frac{\partial^2}{\partial y^2}\right)v = 0,
\end{equation}

and by (2.4),

\begin{equation}\tag{2.7}
\|v\|_{[-A/2,\,A/2]\times B_r(x_1)} \le c\psi(r)\|u\|_{B_r(x_2)}A^{1/2}.
\end{equation}

As explained in \S0, we can now combine the three facts emphasized in [10]
((0.3), $S_D=\Phi(D)$, and (0.6)), to obtain the estimate

\begin{equation}\tag{2.8}
\left|k_f(u)\right|_{B_{A/2}(x_1)}
\le k_nA^{-n/2}\bigl[\tilde{\omega}(B_A(x_1))\bigr]^{-(n+1)/2}\psi(r)\|u\|_{B_A(x_1)}.
\end{equation}

This implies that for all $x_1$,

\begin{equation}\tag{2.9}
\|k_f(x_1,x_2)\|_{B_A(x_2)}
\le k_nA^{-n/2}\bigl[\tilde{\omega}(B_A(x_1))\bigr]^{-(n+1)/2}\psi(r),
\end{equation}

where $k_n$ is a constant depending only on $n$. Similarly, we obtain

\begin{equation}\tag{2.10}
\|\Delta_2 k_f(x_1,x_2)\|_{B_A(x_2)}
\le k_nA^{-n/2}\bigl[\tilde{\omega}(B_A(x_1))\bigr]^{-(n+1)/2}
\left(\frac{2l}{2A}\right)^{2l}\psi(r).
\end{equation}

Then for each fixed $x_1$ we can set

\begin{equation}\tag{2.11}
w(x_1,x_2,y)=\sum_{l=0}^{\infty}\frac{y^{2l}}{(2l)!}\Delta_2^l f(\sqrt{-\Delta})\delta_{x_2},
\end{equation}

to get a harmonic function on $[-A/2,A/2]\times B_A(x_2)$ whose $L^2$ norm satisfies

\begin{equation}\tag{2.12}
\|w(x_1,x_2,y)\|_{[-A/2,A/2]\times B_A(x_2)}
\le c\|k_f(x_1,x_2)\|_{B_A(x_2)}.
\end{equation}

Finishing the argument as above gives

\begin{equation}\tag{2.13}
\left|k_f(x_1,x_2)\right|_{B_{A/2}(x_1)\times B_{A/2}(x_2)}
\le c_n^2A^{-n}\bigl[\tilde{\omega}(B_A(x_1))\tilde{\omega}(B_A(x_2))\bigr]^{-(n+1)/2}\psi^2(r)
\end{equation}

for $\rho(x_1,x_2)=d=2A+r$. The same argument shows that for all $x_1,x_2$,

\begin{equation}\tag{2.14}
\left|k_f(x_1,x_2)\right|_{B_{A/2}(x_1)\times B_{A/2}(x_2)}
\le c_n^2A^{-n}\bigl[\tilde{\omega}(B_A(x_1))\tilde{\omega}(B_A(x_2))\bigr]^{-(n+1)/2}\psi^2(0),
\end{equation}

provided $\varphi\in L^1(\R^+)$. Note that if $k_{f_\lambda}$ denotes the kernel of
$f(\sqrt{-\Delta})$ with respect to the metric $g$, then
$k^g_{f(\lambda)}=c^{-n/2}_g k_{f(\lambda/c)}$. Thus using (2.13), (2.14) we
can estimate $k^g_{f(\lambda)}$ in two different ways. It is straightforward to
check that the estimates so obtained coincide; i.e., (2.13), (2.14) have the
correct scaling property.

As mentioned in \S0, in Proposition 4.2 we give a lower estimate for
$\tilde{\omega}(A,x_0)$, if $A$ is sufficiently small. To simplify the statement
of Theorem 2.1, we will use only Proposition 4.2(i) even though (ii) is sharper
if $\rho(p,x_j)=\rho_j$ is sufficiently large. In \S4 we will define a quantity
$r(H,V)$ by

\begin{equation}\tag{2.15}
V_r^H(H,V)=\frac{V}{2},
\end{equation}

where $V_r^H$ is given by (1.27).

\begin{theorem}[Kernel estimates with Ricci lower bound]
\label{thm:kernel-estimates-ricci-lower-bound}
\dcref{Theorem 2.1}
\group{section-2}
\source{cheeger-gromov-taylor:page-0012}
\textit{Let $M^n$ be a complete Riemannian manifold with $\partial M^n=\emptyset$, and $\operatorname{Ric}_M \ge (n-1)H$. Set $V_{r_0}=V_{r_0}(p)$. If $f \in \mathcal{F}^{\varphi,A}_2$ and $A < \frac12 r(H,V)$, then the kernel $k_f(x_1,x_2)$ is continuous for $x_1 \neq x_2$ and satisfies (2.13), (2.14), where}
\begin{equation}\tag{2.16}
\widetilde{\omega}(B_A(x_j)) \ge \frac{1}{\left(2V_{r+r_0}^{H}/V_{r_0}\right)-1}.
\end{equation}
\end{theorem}

\begin{remark}
\label{rem:gradient-kernel-estimate}
\dcref{Remark after Theorem 2.1}
\group{section-2}
\source{cheeger-gromov-taylor:page-0012}
We point out that just as in [10, \S 4] (or more or less equivalently, by the method of ``domination of semigroups''), we can in fact obtain pointwise bounds on the norm of the \textit{gradient} of $k_f$ under the assumptions of Theorem 2.1. We omit the details.
\end{remark}

In the appendix to this section, we describe sufficient conditions for a function $f$ to be in the class $\mathcal{F}^{\varphi,A}_2$. At this point we will illustrate our results and methods by applying them to the heat kernel.

\begin{example}[The heat kernel]
\label{ex:heat-kernel}
\dcref{Example 2.1}
\group{section-2}
\source{cheeger-gromov-taylor:page-0012,cheeger-gromov-taylor:page-0013,cheeger-gromov-taylor:page-0014,cheeger-gromov-taylor:page-0015}
In this case (2.1)--(2.12) can by bypassed. We begin by recalling the obvious estimate

\begin{equation}\tag{2.17}
\int_r^\infty \frac{e^{-\frac14 \lambda^2/t}}{(4\pi t)^{1/2}}\, d\lambda \le ce^{-\frac14 r^2/t};
\end{equation}

see [26] for a detailed result. Let $u \in L^2$, supp $u \subset B_A(x_2)$. Set $k_{e^{-\lambda^2 t}} = E(x_1,x_2,t)$ and

\begin{equation}\tag{2.18}
v(x_1,t)=\int E(x_1,x_2,t)u(x_2)\, dx_2.
\end{equation}

Assume first that $\rho(x_1,x_2)=d>2A$. By Proposition 1.1 we have for fixed $s$

\begin{equation}\tag{2.19}
\|v(x_1,s)\|_{B_A(x_1)} \le ce^{-\frac14(d-2A)^2/s}\|u\|.
\end{equation}

Using the fact that $e^{-\frac14(d-2A)^2/s}$ is an increasing function of $s$, we have an obvious estimate for $\|v(x_1,s)\|_{B_A(x_1)\times[\frac{z}{2},z]}$. Then using (0.3) we can apply steps 1--3 of \S0 to get

\begin{equation}\tag{2.20}
|v(x_1,z)| \le c_1 c_n \left[\frac{1}{z^{(n+2)/2}} + \frac{1}{A^{n+2}}\right]^{1/2} e^{-\frac14(d-2A)^2/z}
\times \left(\widetilde{\omega}(B_A(x_1))\right)^{-\frac12(n+1)}\|u\|.
\end{equation}

Thus for fixed $x_1, z$,

\begin{equation}\tag{2.21}
\|E(x_1,x_2,z)\| \le c_1 c_n \left[\frac{1}{z^{n/2}} + \frac{z}{A^{n+2}}\right] e^{-\frac14(d-2A)^2/z}
\times \left(\widetilde{\omega}(B_A(x_1))\right)^{-\frac12(n+1)},
\end{equation}

and applying steps 1--3 on $B_A(x_2)\times [\frac{t}{2},t]$, we have

\begin{equation}\tag{2.22}
\begin{aligned}
|E(x_1,x_2,t)|
&\le c_2 c_n^2
\left[\frac{1}{t^{n/2-1}}+\frac{2}{A^{n+2}}\right]
\left[\frac{1}{t^{(n+2)/2}}+\frac{1}{A^{n+2}}\right] \\
&\quad\times e^{-\frac{1}{4}(d-2A)^2/t}
[\tilde\omega(B_A(x_1))\tilde\omega(B_A(x_2))]^{-\frac{1}{2}(n+1)}.
\end{aligned}
\end{equation}

This gives, for $d>2A$,

\begin{equation}\tag{2.23}
|E(x_1,x_2,t)|
\le c_3 c_n^2
\left[\frac{1}{t^{n/2}}+\frac{1}{A^n}\right]
e^{-\frac{1}{4}(d-2A)^2/t}
\times (\tilde\omega(B_A(x_1))\cdot \tilde\omega(B_A(x_2)))^{-\frac{1}{2}(n+1)}.
\end{equation}

Similarly, for all $x_1,x_2$, including those for which $d<2A$,

\begin{equation}\tag{2.24}
|E(x_1,x_2,t)|
\le c_3 c_n^2
\left[\frac{1}{t^{n/2}}+\frac{1}{A^n}\right]
[\tilde\omega(B_A(x_1))\tilde\omega(B_A(x_2))]^{-\frac{1}{2}(n+1)}.
\end{equation}

If we combine (2.23), (2.24) with the estimate of $\tilde\omega(B_A(x))$ in Proposition
4.2, we get a generalization of the estimate of [10]. In this connection we wish
to emphasize that for small time, say $0<t\le 1$, the uniform estimate provided
by (2.24) together with Proposition 4.2 is not a consequence of the usual
asymptotic expansion

\begin{equation}\tag{2.25}
E(x,x,t)\sim \frac{1}{(4\pi t)^{n/2}}\left[1-\frac{1}{6}\tau(x)t+\cdots\right],
\end{equation}

where $\tau$ denotes the scalar curvature. Although we have in fact assumed a
lower bound on $\mathrm{Ric}_{B_A(x)}$, rather than on $\tau(x)$, in effect the upper bound on
$\tau(x)$ has been replaced by a lower bound on $V(B_r(x))$ for some $r$. Scaling the
metric on a flat torus makes it obvious that a lower bound on $V(B_r(x))$ is
really necessary and that (2.25) need not hold uniformly.

As a second example, consider the family of metrics corresponding to a
sharply rounded cone tip $x$, as in Fig. 2.1 below.

\begin{center}
Fig. 2.1
\end{center}

The curvature is $\sim 1/\eps^2$ at the tip and is \emph{nonnegative} everywhere. The volume
$V(B_1(x))$ is bounded below independent of $\eps$. Then according to (2.24),
$E(x,x,t)$ is bounded independent of $\eps$, for $0<t<1$. Again, (2.25) does not
hold uniformly in $\eps$.

We can also derive (2.22) as a consequence of Theorem 2.1. This requires
estimating $\fhat{t}(s)$ and its derivatives, where

\begin{equation*}
(2.26)\qquad f_t(\lambda)=e^{-t\lambda^2}.
\end{equation*}

Note that

\begin{equation*}
(2.27)\qquad \fhat{t}(s)=t^{-1/2}\fhat{1}(t^{-1/2}s),
\end{equation*}

which implies

\begin{equation*}
(2.28)\qquad \fhatk{t}{(k)}(s)=t^{-\frac12(k+1)}\fhatk{1}{(k)}(t^{-1/2}s).
\end{equation*}

We will estimate $\fhatk{1/4}{(k)}(s)$, since

\begin{equation*}
(2.29)\qquad \fhat{1/4}(s)=\pi^{-1/2}e^{-s^2}.
\end{equation*}

Then

\begin{equation*}
(2.30)\qquad \fhatk{1/4}{(k)}(s)=(-1)^k\Hk(s)e^{-s^2},
\end{equation*}

where $H_k(s)$ is the $k$ th Hermite polynomial; see Lebedev [26, p. 60]. In order
to get a good estimate on $H_k(s)$, it is convenient to use the generating function
identity

\begin{equation*}
(2.31)\qquad \sum_{n=0}^{\infty}\frac{\Hn(s)^2}{2^n n!}\eta^n=(1-\eta^2)^{-1/2}e^{2s^2\eta/(1+\eta)}, \qquad |\eta|<1;
\end{equation*}

see [26, p. 61]. Pick $\eta$ fairly small, so that

\begin{equation*}
(2.32)\qquad \frac{|\Hn(s)|^2}{2^n n!}|\eta|^n \le c e^{\eps s^2},
\end{equation*}

with $\eps>0$ given. It follows that

\begin{equation*}
(2.33)\qquad \left|\fhatk{1/4}{(k)}(s)\right|\le c\left[\left(\frac{2}{|\eta|}\right)^k k!\right]^{1/2}e^{-(1-\eps)s^2}.
\end{equation*}

Consequently,

\begin{equation*}
(2.34)\qquad
\begin{aligned}
\left|\fhatk{t}{(k)}(s)\right|
&\le c\left[\left(\frac{2}{|\eta|}\right)^k k!\right]^{1/2}t^{-\frac12(k+1)}e^{-\frac14(1-\eps)s^2/t}\\
&= c\left[\left(\frac{2}{|\eta|t}\right)^k k!\right]^{1/2}t^{-1/2}e^{-\frac14(1-\eps)s^2/t}.
\end{aligned}
\end{equation*}

Thus we can take $f_t \in \mathscr{F}_2^{\varphi,A}$ with
\[
\varphi(r)=t^{-1/2}e^{-\frac14(1-\varepsilon)r^2/t},
\]
and, by Theorem 2.1, recover essentially the same estimates (2.24), (2.25) on the heat kernel.
\end{example}

To close this section, we show how, for manifolds of dimension $2$ or $3$, the more general class of estimates of \S1 can in fact be derived under the weaker assumption that Ricci curvature is bounded from below.

Let
\begin{equation}\tag{2.35}
\Hg=e_i e_j(g)-\nabla_{e_i}\nabla_{e_j}(g)
\end{equation}
denote the Hessian of $g$. We will show that the constant $c(r)$ in the coercive estimate
\begin{equation}\tag{2.36}
\|Hg\|_{B_{r/2}(x_0)}^{2}\le c(r)\Bigl(\|dg\|_{B_r(x_0)}^{2}+\|\Delta g\|_{B_r(x_0)}^{2}\Bigr)
\end{equation}
can be bounded from above in terms of $r$, $H$, if we assume $\Ric_M\ge H$; compare [10, \S4].

However, the usual proof of the Sobolev inequality which estimates pointwise norm in terms of $H^s$, $s>n/2$, norm shows that the constant in that inequality can be bounded from above if a lower bound for the refined injectivity angle $\widetilde{\omega}(H,\varepsilon,r_0,x_0)$, (as defined before Proposition 2.3), is known.
If we combine this with (2.36), we find that if dim $M=2,3$, we obtain directly a pointwise estimate on $k_f$, $f\in \mathscr{F}_1^{1}$, only assuming $\Ric_M\ge (n-1)H$, $V_{r_0}(p)\ge V$. We turn to the proof of (2.36).

Recall that for any 1-form $\omega$ we have the Bochner-Weitzenbock formula
\begin{equation}\tag{2.37}
(d\delta+\delta d)\omega=\sum_i\bigl(-\nabla_{e_i}\nabla_{e_i}+\nabla_{\nabla_{e_i}e_i}\bigr)\omega+\Ric(\omega^*),
\end{equation}
where
\begin{equation}\tag{2.38}
\Ric(\omega^*)=(\Ric(\omega))^*,
\end{equation}
and $\omega^*$ is the tangent vector dual to $\omega$. Fix a metric ball $B_r(x)$, and let $\psi(r)$ be defined by $\psi(r)|_{[0,r/2]}\equiv 1$, $\psi(r)|_{[r/2,r]}\equiv -2r/\tilde r+2$. By Stokes' theorem,
\begin{equation}\tag{2.39}
\int_{B_{\tilde r}(x_0)}\psi^2\, d\delta dg\wedge *dg
=\int_{B_{\tilde r}(x_0)}\psi^2\delta dg\wedge *\delta dg
-\int_{B_{\tilde r}(x_0)}2\psi\, d\psi\delta dg\wedge *dg.
\end{equation}

Applying (2.37) to the left-hand side and integrating by parts gives
\[
\int_{B_r(x_0)} \psi^2 \delta \dg \wedge *\dg
= \int_{B_r(x_0)} \left\{ \psi^2 \sum_i \left(-\nabla_{e_i}\nabla_{e_i} + \nabla_{\nabla_{e_i} e_i}\right)\dg \wedge *\dg + \psi^2 \Ric(\dg) \wedge *\dg \right\}
\]
\begin{equation}
= \int_{B_r(x_0)} \left\{ \psi^2 \sum_i \nabla_{e_i}\dg \wedge *\nabla_{e_i}\dg + 2\psi \sum_i e_i(\psi)\nabla_{e_i}\dg \wedge *\dg
+ \psi^2 \Ric(\dg) \wedge *\dg \right\}.
\tag{2.40}
\end{equation}

If we substitute (2.40) into (2.39) and transpose the middle term we get
\[
\int_{B_r(x_0)} \left( \psi^2 \sum_i \nabla_{e_i}\dg \wedge *\nabla_{e_i}\dg + \Ric(\dg)\wedge *\dg \right)
\]
\begin{equation}
= \int_{B_r(x_0)} \psi^2 \delta \dg \wedge *\delta \dg
- \int_{B_r(x_0)} 2\psi\, d\psi\,\delta \dg \wedge *\dg
- \int_{B_r(x_0)} 2\psi \sum_i e_i(\psi)\nabla_{e_i}\dg \wedge *\dg.
\tag{2.41}
\end{equation}

Note that $\nabla \dg = Hg$. Applying the Schwarz inequality to the last two terms of the right-hand side of (2.41) gives
\[
\int_{B_r} \psi^2 \sum_i \nabla_{e_i}\dg \wedge *\nabla_{e_i}\dg + (n-1)H \|\psi\,dg\|^2_{B_r(x_0)}
\]
\begin{equation}
\begin{aligned}
&\le \|\Delta g\|^2_{B_r(x_0)}
+ \left(\|\Delta g\|_{B_r(x_0)} + \frac{4}{r^2}\|dg\|^2_{B_r(x_0)}\right) \\
&\quad + \left(\frac12\int_{B_r(x_0)} \psi^2 \sum_i \nabla_{e_i}\dg \wedge *\nabla_{e_i}\dg
+ \frac{8}{r^2}\|dg\|^2_{B_r(x_0)}\right).
\end{aligned}
\tag{2.42}
\end{equation}

If we transpose the fourth term on the right-hand side, we obtain

\begin{proposition}[Coercive Hessian estimate]
\label{prop:coercive-hessian-estimate}
\dcref{Proposition 2.2}
\group{section-2}
\source{cheeger-gromov-taylor:page-0015,cheeger-gromov-taylor:page-0016}
\textit{Let $M^n$ be complete, and $\Ric_M \geq (m-1)H$. Then for all $x, r$}
\begin{equation}\tag{2.43}
\|Hg\|^2_{B_{r/2}(x)} \le 4\|\Delta g\|^2_{B_r(x)} + \frac{24}{r^2}\|dg\|^2_{B_r(x)}
- 2(n-1)H\|\psi\,dg\|^2_{B_r(x)}.
\end{equation}
\end{proposition}

Let $A(r,\theta), \mathcal{Q}^H(r)$ be as in \S1; see (1.28). Given $\epsilon > 0$ and $H$, we define a refinement of the injectivity angle by letting $\widetilde{\omega}(H,\epsilon,r,x)$ be the angular
measure of the set $\Omega^H(\epsilon)$ of initial tangent vectors to geodesics $\gamma$ from $x$ such
that $\gamma|_{[0,r]}$ is minimal and

\begin{equation}\tag{2.44}
\frac{A(s,\theta)}{\mathcal{Q}^H(s)} \ge \epsilon,\qquad 0 \le s \le r,
\end{equation}

at $\gamma(s)$. Note that for all $H$, $\widetilde{\omega}(H,0,r,x)=\widetilde{\omega}(r,x)$. Let $\nabla^i g$ denote the $i$-th
covariant derivative of $g$.

\begin{proposition}[Pointwise Sobolev estimate]
\label{prop:pointwise-sobolev-estimate}
\dcref{Proposition 2.3}
\group{section-2}
\source{cheeger-gromov-taylor:page-0017}
\textit{Let $g$ be a smooth function on $B_r(x_0)\subset M^n$. Then for}
$N>n/2$,
\begin{equation}\tag{2.45}
g^2(x_0)\le \frac{c_3(N)}{\epsilon\,\widetilde{\omega}(H,\epsilon,\bar r,x_0)}
\left(\int_0^{\bar r} r^{2(N-1)}\,\frac{dr}{\mathcal{Q}^H(r)}\right)
\sum_{i=0}^N \bar r^{2(i-N)}\|\nabla^i g\|^2_{B_{\bar r}(x_0)}.
\end{equation}
\end{proposition}

\begin{proof}
\source{cheeger-gromov-taylor:page-0017}
Let $\psi$ be essentially as above, but made $C^\infty$. Then
\begin{equation}
g^2(x_0)\le \frac{1}{\widetilde{\omega}(H,\epsilon,\bar r,x_0)}
\int_{\Omega^H(\epsilon)}\left|\int_0^{\bar r}\frac{d}{dr}\bigl(\psi g(r,\theta)\bigr)\,dr\right|^2 d\theta
\end{equation}
\begin{equation}\tag{2.46}
= \frac{1}{\widetilde{\omega}(H,\epsilon,\bar r,x_0)}
\cdot \int_{\Omega^H(\epsilon)}\left[\int_0^{\bar r}\frac{r^{N-1}}{A^{1/2}(r,\theta)}\frac{d^N}{dr^N}\bigl(\psi g(r,\theta)\bigr)A^{1/2}(r,\theta)\,dr\right]^2 d\theta
\end{equation}
\[
\le \frac{c_3(N)}{\epsilon\,\widetilde{\omega}(H,\epsilon,\bar r,x_0)}
\left(\int_0^{\bar r}\frac{r^{2(N-1)}\,dr}{\mathcal{Q}^H(r)}\right)
\sum_{i=0}^N \bar r^{2(i-N)}\int_{B_{\bar r}(x_0)}\|\nabla^i g\|^2,
\]
and the claim follows.
\end{proof}

In Proposition 4.2 we will give a lower estimate for $\widetilde{\omega}(H,\epsilon,\bar r,x_0)$ for $\epsilon,\bar r$
sufficiently small. We will now combine this with Propositions 2.2 and 2.3 to
get the following result. As in previous instances, the statement can be
sharpened somewhat; see Proposition 4.2(iv).

\begin{theorem}[Ricci-lower-bound pointwise estimate]
\label{thm:ricci-lower-bound-pointwise-estimate}
\dcref{Theorem 2.4}
\group{section-2}
\source{cheeger-gromov-taylor:page-0017,cheeger-gromov-taylor:page-0018}
\textit{Let $M^n$ be complete, and $\mathrm{Ric}_M\ge (n-1)H$. Fix $p, x\in M^n$}
\textit{and $r_0>0$. Set $V_{r_0}=V_{r_0}(p)$, $\rho=\rho(p,x)$, $V^H_{r(H,V_{r_0})}=\frac12 V_{r_0}$.}

\noindent (i) Then for $N>n/2$,
\begin{equation}\tag{2.47}
g^2(x)\le \frac{c_3(N)}{\bigl[(4V^H_{\rho+r_0}/V_{r_0})-3\bigr]^2}
\left(\int_0^{r(H,V_{r_0})}\frac{r^{2(N-1)}\,dr}{\mathcal{Q}^H(r)}\right)
\sum_{i=0}^N \bar r^{2(i-N)}\|\nabla^i g\|^2_{B_r(x)}.
\end{equation}

\noindent (ii) In particular, for $n=2,3$,
\begin{equation}\tag{2.48}
\begin{aligned}
g^2(x)&\le \frac{c_3(2)}{\bigl[(4V^H_{\rho}/V_{r_0})-3\bigr]^2}
\left(\int_0^{r(H,V_{r_0})}\frac{r^2}{\mathcal{Q}^H(r)}\,dr\right) \\
&\quad\times\Bigl[r_0^{-4}\|g\|^2_{B_{r_0}(x)}
+25r_0^{-2}\|dg\|^2_{B_{r_0}(x_0)} \\
&\qquad -2(n-1)H\|\psi dg\|^2_{B_{r_0}(x_0)}
+4\|\Delta g\|^2_{B_{r_0}(x_0)}\Bigr].
\end{aligned}
\end{equation}
\end{theorem}

In view of (2.48) the kernels $k_f$, $f\in \mathscr{F}_1^1$, can now be treated by the method of \S1.

\begin{center}
\textbf{Appendix to \S2. The class $\mathscr{F}_2^{\varphi,A}$, an example}
\end{center}

As we saw above, the functions
\begin{equation}\tag{A.1}
f_t(\lambda)=e^{-t\lambda^2}
\end{equation}
are elements of $\mathscr{F}_2^{\varphi,A}$ with
\begin{equation}\tag{A.2}
\varphi(r)=t^{-1/2}e^{-\frac14 r^2/t},\qquad t'>t.
\end{equation}
Here we want to give another example of a large class of functions which
belong to $\mathscr{F}_2^{\varphi,A}$ with
\begin{equation}\tag{A.3}
\varphi(r)=e^{-Wr}.
\end{equation}
Consider a region in the complex domain of the form
\begin{equation}\tag{A.4}
\Omega=\{z\in \mathbf{C}: |\operatorname{Im} z|\le W\}\cup\{z\in \mathbf{C}: |\operatorname{Im} z|\le B\,|\operatorname{Re} z|\}.
\end{equation}
Suppose $f$ is holomorphic in $\Omega$ and satisfies the estimate
\begin{equation}\tag{A.5}
|f(z)|\le c(1+|z|)^m,\qquad z\in \Omega.
\end{equation}
If $W$ and $B$ are decreased slightly, it follows from Cauchy's integral formula
that
\begin{equation}\tag{A.6}
|f^{(k)}(z)|\le c(ck)^k(1+|z|)^{m-k},
\end{equation}
and in particular,
\begin{equation}\tag{A.7}
|z^k f^{(k+m)}(z)|\le c(ck)^k.
\end{equation}
Now if $g(z)$ is holomorphic on $|\operatorname{Im} z|\le W$, then
\begin{equation}\tag{A.8}
\int_{-\infty}^{\infty}|g(x+iy)|\,dx\le G,\qquad |y|\le W,
\end{equation}
implies
\begin{equation}\tag{A.9}
|\hat g(s)|\le Ge^{-W|s|},\qquad s\in \mathbf{R}.
\end{equation}

Thus from (A.7) it follows that
\[
\tag{A.10}
\lvert s^{k+m}\hat{f}^{(k)}(s)\rvert \le c_1 (ck)^k e^{-W\lvert s\rvert},
\]
so
\[
\tag{A.11}
\lvert \hat{f}^{(k)}(s)\rvert \le c_1 (ck)^k e^{-W\lvert s\rvert}, \qquad \lvert s\rvert \ge 1.
\]
Thus $f \in \mathcal{G}^{\varphi,A}_{2}$ with $\varphi(r)$ given by (A.3).
